\section{Conservation Laws and Shocks}\label{chap:conservation-laws-shocks}\label{chap:local-existence}

Recall that in Problem~\ref{ex:discontinuous-weak-wave-solutions} we defined a weak solution of the one-dimensional wave equation to be a function $u(x,t)$ such that
\begin{equation}\label{eq:conservation-intro-weak-wave}
\int_{-\infty}^{\infty}\int_{-\infty}^{\infty}u(x,t)(\phi_{tt}(x,t)-\phi_{xx}(x,t))\,\mathrm{d}x\,\mathrm{d}t=0
\end{equation}
for every $\phi\in C_0^2(\mathbb{R}^2)$. In the problem, one was asked to show the following:
\begin{enumerate}
\item that any strong (classical $C^2$) solution of the wave equation is also a weak solution,
\item that discontinuous functions of the form
\begin{equation}\label{eq:conservation-intro-right-step}
u(x,t):=H(x-t),
\end{equation}
and
\begin{equation}\label{eq:conservation-intro-left-step}
u(x,t):=H(x+t),
\end{equation}
where $H$ is the Heaviside function,
\begin{equation}\label{eq:conservation-intro-heaviside}
H(x):=\begin{cases}0,&x<0\\1,&x\geq0\end{cases}
\end{equation}
are weak solutions of the wave equation.\end{enumerate}

In this section, we extend the notion of weak solution to strictly hyperbolic systems of conservation laws in one space dimension. As in the case of the wave equation, we will be able to observe discontinuous solutions. The use of the term \emph{shock wave} to describe a discontinuous weak solution originated in the study of gas dynamics, but now it is used in connection with any discontinuous solution of a hyperbolic problem.

\subsection{Systems in One Space Dimension}\label{sec:one-dimensional-conservation-systems}

We consider the Cauchy problem for the following first-order quasilinear system of equations
\begin{equation}\label{eq:conservation-law-system}
\bm{u}_t+\bm{f}(\bm{u})_x=0.
\end{equation}
Here $(x,t)$ lies in the upper half-plane $\mathbb{R}^{2+}:=(-\infty,\infty)\times[0,\infty)$; the unknown is $\bm{u}:\mathbb{R}^{2+}\to\mathbb{R}^n$, and $\bm{f}:\mathbb{R}^n\to\mathbb{R}^n$ is a given function assumed to be sufficiently smooth (usually at least $C^2(D)$, for some open set $D\subseteq\mathbb{R}^n$). We refer to the space $\mathbb{R}^n$ in which the dependent variables of the unknown lie as \emph{state space}. Our initial condition is
\begin{equation}\label{eq:conservation-law-cauchy-data}
\bm{u}(x,0)=\bm{u}_0(x)
\end{equation}
for $x\in\mathbb{R}$, where the function $\bm{u}_0:\mathbb{R}\to\mathbb{R}^n$ is given.

The following examples will be mentioned repeatedly throughout this chapter.

\begin{Example}{Burgers' equation}
\label{ex:burgers-equation}
The following equation is the most commonly used example of a single $(n=1)$ nonlinear conservation law.
\begin{equation}\label{eq:burgers-conservation-law}
u_t+\left(\dfrac{u^2}{2}\right)_x=u_t+uu_x=0.
\end{equation}
Here $u:\mathbb{R}^{2+}\to\mathbb{R}$. The equation was studied by Burgers and Hopf, and was considered to be a crude model for turbulence.
\end{Example}

\begin{Example}{The wave equation}
\label{ex:wave-first-order-system}
The linear one-dimensional wave equation \eqref{eq:wave-one-dimensional} can be converted from a second-order single equation into a first-order system by setting $v=u_t$ and $w=u_x$. The resulting system is
\begin{equation}\label{eq:wave-first-order-w}
w_t=v_x,
\end{equation}

\begin{equation}\label{eq:wave-first-order-v}
v_t=w_x.
\end{equation}
Of course, the first equation is simply the equality of mixed partials of $u$, whereas the second is the original wave equation.
\end{Example}

\begin{Example}{The $p$ system}
\label{ex:p-system}
This is simply the first-order system derived from the nonlinear wave equation
\[
u_{tt}=\phi(u_x)_x,
\]
where $\phi:\mathbb{R}\to\mathbb{R}$ is a given function. Setting $v=u_t$ and $w=u_x$ as we did for the wave equation yields
\begin{equation}\label{eq:nonlinear-wave-first-order-w}
w_t=v_x,
\end{equation}

\begin{equation}\label{eq:nonlinear-wave-first-order-v}
v_t=\phi(w)_x.
\end{equation}
It is very common to define $p(w):=-\phi(w)$ and write the system in the form
\begin{equation}\label{eq:p-system-volume}
w_t-v_x=0,
\end{equation}

\begin{equation}\label{eq:p-system-velocity}
v_t+p(w)_x=0.
\end{equation}
A common model for isentropic gas dynamics uses this system with
\begin{equation}\label{eq:isentropic-pressure-law}
p(w):=kw^{-\gamma},
\end{equation}
and with the restriction $w>0$. Here $\gamma>1$ and $k>0$ are constants.
\end{Example}

\begin{Example}{Gas dynamics in Lagrangian coordinates}
\label{ex:lagrangian-gas-dynamics}
The following equations describe the motion of an inviscid gas that does not conduct heat:
\begin{equation}\label{eq:lagrangian-gas-mass}
v_t-u_x=0,
\end{equation}

\begin{equation}\label{eq:lagrangian-gas-momentum}
u_t+p_x=0,
\end{equation}

\begin{equation}\label{eq:lagrangian-gas-energy}
E_t+(up)_x=0.
\end{equation}
Here $v$ is the \emph{specific volume}, $u$ is the \emph{velocity}, $p$ is the \emph{pressure} and $E$ is the \emph{specific energy per unit mass}. The specific volume is defined to be the reciprocal of the \emph{density} $\rho$
\begin{equation}\label{eq:specific-volume-density}
v:=1/\rho.
\end{equation}
Equation \eqref{eq:lagrangian-gas-mass} represents conservation of mass, \eqref{eq:lagrangian-gas-momentum} represents conservation of linear momentum and \eqref{eq:lagrangian-gas-energy} represents conservation of energy.

In order to make this system of three equations in four unknowns well-posed, we must add a \emph{constitutive equation} or \emph{equation of state} that describes one of the variables as a given function of the other three. This is done here with the pressure, which is usually given by
\begin{equation}\label{eq:gas-pressure-state-equation}
p:=\hat p(e,v),
\end{equation}
where $e:=E-u^2/2$ is the \emph{internal energy}.
\end{Example}

\begin{Example}{Gas dynamics in Eulerian coordinates}
\label{ex:eulerian-gas-dynamics}
In the \emph{Lagrangian} description of gas dynamics above, the variable $x$ describes a fixed particle of gas. In the \emph{Eulerian} description, $x$ describes a fixed point in space. When the equations are derived using such a model, the following system of equations results:
\begin{equation}\label{eq:eulerian-gas-mass}
\rho_t+(\rho u)_x=0,
\end{equation}

\begin{equation}\label{eq:eulerian-gas-momentum}
(\rho u)_t+(\rho u^2+p)_x=0,
\end{equation}

\begin{equation}\label{eq:eulerian-gas-energy}
\left[\rho\left(\dfrac{u^2}{2}+e\right)\right]_t+\left[\rho u\left(\dfrac{u^2}{2}+i\right)\right]_x=0.
\end{equation}
Here $\rho$, $u$, $p$ and $e$ are defined as above; and $i:=e+p/\rho$ is the \emph{specific enthalpy}. Similarly to the equations above, \eqref{eq:eulerian-gas-mass} represents conservation of mass, \eqref{eq:eulerian-gas-momentum} represents conservation of linear momentum and \eqref{eq:eulerian-gas-energy} represents conservation of energy.
\end{Example}
